\documentclass[10pt,journal,compsoc]{IEEEtran}

\usepackage[utf8]{inputenc}
\usepackage{amsmath,amssymb,amsfonts}
\usepackage{graphicx}
\usepackage{booktabs}
\usepackage{hyperref}
\usepackage{subcaption}
\usepackage{microtype}
\usepackage{cite}

\hypersetup{
    colorlinks=true,
    linkcolor=blue,
    citecolor=cyan,
    urlcolor=blue
}

\begin{document}

\title{Energy-Based Models and Score-Based Generative Modeling with FiLM\\
\large Assignment 3 -- Deep Generative Models (Fall 1404 / 2025)}

\author{Alireza Najafi Motiei (Student ID: 810100224)\\
\IEEEmembership{Faculty of Electrical and Computer Engineering, University of Tehran}\\
Instructor: Dr. Mostafa Tavassolipour}

\markboth{University of Tehran, ECE Department -- Deep Generative Models (Fall 1404)}{Najafi Motiei: Assignment 3 Report}

\maketitle

\begin{abstract}
This report presents the theoretical derivations and experimental results for Assignment 3 of Deep Generative Models (Fall 1404). The first module explores Energy-Based Models (EBMs), analyzing the intractability of the partition function, the maximum likelihood gradient dynamics, contrastive divergence (CD), and the exponential failure of rejection sampling in high dimensions ($D=784$). We implement an EBM on MNIST and evaluate Langevin dynamics for image synthesis and denoising across noise variances ($\sigma \in \{0.2, 0.4, 0.6\}$). The second module investigates Score-Based Generative Models (SGMs). We examine score matching and denoising score matching (DSM) to bypass Jacobian trace computations, analyze mode-mixing pathology and manifold hypothesis challenges, and demonstrate Noise Conditional Score Networks (NCSN) combined with Annealed Langevin Dynamics and Feature-wise Linear Modulation (FiLM) conditioning.
\end{abstract}

\begin{IEEEkeywords}
Energy-Based Models, Partition Function, Contrastive Divergence, Langevin Dynamics, Score-Based Generative Models, Denoising Score Matching, Annealed Langevin Dynamics, FiLM Conditioning, MNIST.
\end{IEEEkeywords}

\section{Question 1: Energy-Based Models (EBM)}

\subsection{Theoretical Foundations and the Partition Function}
An Energy-Based Model parameterizes a probability distribution as:
\begin{equation}
p_\theta(x) = \frac{\exp(-E_\theta(x))}{Z(\theta)}, \quad Z(\theta) = \int \exp(-E_\theta(x)) dx
\end{equation}
Computing the normalizing partition function $Z(\theta)$ requires integrating over high-dimensional input spaces, which is computationally intractable.

\subsubsection{Maximum Likelihood Learning Gradient}
The gradient of the log-likelihood decomposes into positive and negative phases:
\begin{equation}
\nabla_\theta \log p_\theta(x) = -\nabla_\theta E_\theta(x) + \mathbb{E}_{x' \sim p_\theta(x')} [\nabla_\theta E_\theta(x')]
\end{equation}
The positive phase suppresses the energy of observed data samples, while the negative phase elevates the energy of model fantasies drawn from the current distribution.

\subsubsection{Contrastive Divergence (CD)}
Evaluating the expectation $\mathbb{E}_{p_\theta}$ requires Markov Chain Monte Carlo (MCMC). Running chains to equilibrium at every gradient step is infeasible. Contrastive Divergence ($CD_k$) addresses this by initializing the Markov chain at data points $x$ and executing only a small number of $k$ steps (e.g., $k=1$ or $k=10$), effectively penalizing model perturbations away from the data manifold.

\subsubsection{Curse of Dimensionality in Rejection Sampling}
Consider generating MNIST digits ($D = 784$) via rejection sampling with proposal $q(x) = \mathcal{N}(0, \sigma_q^2 I)$ and target $p(x) = \mathcal{N}(0, \sigma_p^2 I)$, where $\sigma_q = 1.01 \sigma_p$ (a mere $1\%$ variance discrepancy). The acceptance probability is upper-bounded by:
\begin{equation}
P(\text{accept}) = \frac{1}{M} \approx \left(\frac{\sigma_p}{\sigma_q}\right)^D = \left(\frac{1}{1.01}\right)^{784} \approx 4.02 \times 10^{-4}
\end{equation}
On average, roughly $2,500$ proposal evaluations are required to accept a single valid sample, rendering rejection sampling completely impractical for high-dimensional generative modeling.

\subsection{Experimental Implementation on MNIST}
We trained an EBM parameterized by a deep convolutional network for 15 epochs utilizing ReduceLROnPlateau scheduling and Langevin sampling.

\begin{figure}[htbp]
\centering
\includegraphics[width=0.95\columnwidth]{figures/image4.png}
\caption{EBM training progression: Loss metrics, positive/negative energy curves ($E_{\text{real}}, E_{\text{fake}}$), and Langevin samples.}
\label{fig:ebm_loss}
\end{figure}

\begin{figure}[htbp]
\centering
\includegraphics[width=0.85\columnwidth]{figures/image6.png}
\caption{Synthesized digit samples generated from uniform random noise via Langevin dynamics.}
\label{fig:ebm_samples}
\end{figure}

\subsubsection{Image Denoising Across Varied Noise Levels}
Langevin dynamics was evaluated for image restoration across noise standard deviations $\sigma \in \{0.2, 0.4, 0.6\}$.
\begin{itemize}
    \item \textbf{$\sigma = 0.2$}: The model successfully removes background artifacts and restores sharp digit strokes.
    \item \textbf{$\sigma = 0.4$}: High recovery of topological loops, with minor edge oversmoothing.
    \item \textbf{$\sigma = 0.6$}: The initial sample is severely corrupted; the EBM produces a clean, recognizable digit, occasionally altering digit identity (e.g., collapsing a corrupted 2 into a 1) to land in the nearest energy well.
\end{itemize}

\begin{figure}[htbp]
\centering
\includegraphics[width=0.9\columnwidth]{figures/image9.png}
\caption{Denoising results at $\sigma = 0.6$ demonstrating background noise removal and topological restoration.}
\label{fig:denoise_06}
\end{figure}

\section{Question 2: Score-Based Generative Modeling}

\subsection{The Score Function and Partition Function Independence}
The score function of a density $p(x)$ is defined as the gradient of its log-density:
\begin{equation}
s(x) \triangleq \nabla_x \log p(x)
\end{equation}
Substituting the unnormalized energy parameterization:
\begin{equation}
s(x) = \nabla_x [ -E_\theta(x) - \log Z(\theta) ] = -\nabla_x E_\theta(x)
\end{equation}
Because $Z(\theta)$ is independent of $x$, its spatial gradient is identically zero. Hence, score matching entirely circumvents the partition function!

\subsection{Denoising Score Matching (DSM)}
Standard score matching minimizes:
\begin{equation}
J_{\text{SM}}(\theta) = \mathbb{E}_{p_{\text{data}}(x)} \left[ \frac{1}{2} \|s_\theta(x)\|_2^2 + \text{tr}(\nabla_x s_\theta(x)) \right]
\end{equation}
Evaluating the Jacobian trace $\text{tr}(\nabla_x s_\theta(x))$ requires $D$ separate backpropagation passes, scaling prohibitively for image generation.

Denoising Score Matching adds Gaussian noise $\tilde{x} \sim \mathcal{N}(x, \sigma^2 I)$ and trains a network to predict the analytic perturbation vector:
\begin{equation}
\nabla_{\tilde{x}} \log q_\sigma(\tilde{x} \mid x) = -\frac{\tilde{x} - x}{\sigma^2}
\end{equation}
The DSM objective simplifies to standard mean squared error:
\begin{equation}
\mathcal{L}_{\text{DSM}}(\theta) = \mathbb{E}_{x \sim p_{\text{data}}} \mathbb{E}_{\tilde{x} \sim \mathcal{N}(x, \sigma^2 I)} \left[ \left\| s_\theta(\tilde{x}) + \frac{\tilde{x} - x}{\sigma^2} \right\|_2^2 \right]
\end{equation}
which requires only a single forward-backward pass per mini-batch.

\subsection{Pathologies of Single-Scale Score Matching}
\begin{enumerate}
    \item \textbf{Manifold Hypothesis}: Data concentrates on low-dimensional manifolds, leaving scores undefined in ambient empty space.
    \item \textbf{Global Weight Blindness}: Scores retain only local derivative slopes, ignoring relative mode heights ($\pi_A$ vs. $\pi_B$).
    \item \textbf{Lack of Mixing}: Isolated modes separated by low-density chasms experience vanishing gradients, trapping Langevin sampling in local basins.
\end{enumerate}

\subsection{NCSN and Annealed Langevin Dynamics with FiLM}
Song \& Ermon resolve these challenges by perturbing data with a geometric noise sequence $\{\sigma_i\}_{i=1}^L$ ($\sigma_1 > \sigma_2 > \dots > \sigma_L$). Large $\sigma_1$ fills low-density valleys, connecting isolated modes and enabling global exploration. Annealing down to small $\sigma_L$ restores fine-grained visual details.

\subsubsection{FiLM Conditioning Mechanism}
To condition the score network on noise levels $\sigma_i$ or class labels $y$, Feature-wise Linear Modulation (FiLM) modulates intermediate feature activations $h$:
\begin{equation}
\text{FiLM}(h \mid \sigma) = \gamma(\sigma) \odot h + \beta(\sigma)
\end{equation}
where affine scales $\gamma$ and shifts $\beta$ are learned projections of the conditioning scalar.

\begin{figure}[htbp]
\centering
\includegraphics[width=0.95\columnwidth]{figures/image22.png}
\caption{Vector fields and simulated Annealed Langevin trajectories navigating from noise toward data clusters.}
\label{fig:score_field}
\end{figure}

\begin{figure}[htbp]
\centering
\includegraphics[width=0.9\columnwidth]{figures/image27.png}
\caption{Sampling convergence demonstrating coarse-to-fine sample generation across the noise schedule.}
\label{fig:score_samples}
\end{figure}

\bibliographystyle{IEEEtran}
\begin{thebibliography}{1}
\bibitem{song2019generative}
Y.~Song and S.~Ermon, ``Generative modeling by estimating gradients of the data distribution,'' in \emph{NeurIPS}, 2019.
\bibitem{vincent2011connection}
P.~Vincent, ``A connection between score matching and denoising autoencoders,'' \emph{Neural Computation}, vol.~23, no.~7, pp. 1661--1674, 2011.
\bibitem{du2019implicit}
Y.~Du and I.~Mordatch, ``Implicit generation and modeling with energy based models,'' in \emph{NeurIPS}, 2019.
\bibitem{perez2018film}
E.~Perez, F.~Strub, H.~de~Vries, V.~Dumoulin, and A.~Courville, ``FiLM: Visual reasoning with a general conditioning layer,'' in \emph{AAAI}, 2018.
\end{thebibliography}

\end{document}
